% ============================================================
%  chapters/ch22-simulated-agency.tex
% ============================================================

\chapter{Simulated Agency}

\chapterepigraph{%
  Consciousness is not the light that illuminates the room.\\
  It is the process by which the room decides\\
  which lights are worth keeping on.
}{Flyxion, \textit{Semantic Infrastructure}}

\sidenote{App.\,L for\\semantic fields\\and simulated\\agency; connects\\to IAD framework}

Part~V ended with the universe as an accessibility landscape. Part~VI begins
with the question that this raises immediately: what is it for a \emph{part}
of the universe to model the \emph{rest} of the universe? What is it for a
physical system to represent, predict, and act?

This is the question of agency. And the central argument of this chapter
is that agency — including conscious agency — is best understood not as a
special substance or emergent property but as a \emph{projection-based
inference process}: a system that constructs sparse internal simulations
of its environment and navigates those simulations using accessibility
gradients.

\section{Projection-Based Consciousness}

The projection framework of Appendix~B and Chapter~16 gives us the
right starting point. An agent $A$ embedded in an environment $E$ has
access not to $E$ directly but to a projection:
\[
  O = \Pi_A(E) \in \mathcal{O}_A
\]
where $\mathcal{O}_A$ is the agent's observation space. The agent's
observation is always a fiber: many environmental states are consistent
with any given observation.

\begin{definition}{Projection-Based Agent}{proj-agent}
  A \emph{projection-based agent} is a tuple $(E, \Pi_A, \hat{E}_A,
  \mathcal{P}_A)$ where:
  \begin{itemize}
    \item $E$ is the environment state space,
    \item $\Pi_A : E \to \mathcal{O}_A$ is the agent's observation
          projection,
    \item $\hat{E}_A$ is the agent's internal model space (simulations),
    \item $\mathcal{P}_A : \mathcal{O}_A \to \Delta(\hat{E}_A)$ is the
          agent's posterior map — assigning a distribution over internal
          models to each observation.
  \end{itemize}
\end{definition}

Consciousness, in this picture, is the process of maintaining and updating
$\mathcal{P}_A$ — the distribution over internal simulations consistent
with current observations. The agent is not directly conscious of $E$;
it is conscious of its \emph{best current simulation} of $E$.

\begin{conceptaside}{The Simulation Hypothesis Inverted}
  The popular ``simulation hypothesis'' asks whether we might be living
  in a simulation run by some external intelligence. The projection-based
  view inverts this: every agent is \emph{already} living in a simulation
  — specifically, the internal model $\hat{E}_A$ that its own cognitive
  processes maintain. The question is not whether we live in a simulation
  but how faithfully our internal simulation tracks the external environment.

  This inversion has practical consequences. Hallucination, delusion,
  and confabulation are not pathological departures from direct perception
  — they are normal features of the simulation process when internal
  models diverge from external states faster than they can be corrected.
  Perception is always inference from projections; the question is how
  well-constrained the inference is.
\end{conceptaside}

\section{Sparse Inference}

A key feature of the projection-based framework is that internal models
are necessarily \emph{sparse}: they cannot represent the full complexity
of $E$, but only the features relevant to the agent's current goals and
constraints.

\begin{definition}{Sparse Internal Model}{sparse-model}
  An internal model $\hat{e} \in \hat{E}_A$ is \emph{$k$-sparse} if it
  specifies at most $k$ degrees of freedom of the environment, with all
  other degrees left implicit (marginalized over). The \emph{sparsity
  ratio} is $k / \dim(E)$.
\end{definition}

For a human agent, $\dim(E)$ is astronomical — every particle in the
observable universe — while $k$ is of order $10^2$ to $10^4$ (the
number of features we are consciously tracking at any moment). The
sparsity ratio is essentially zero. We are always inferring an
astronomically sparse sketch of reality.

CLIO (Chapter~15) explains why this works: constraint structure allows
sparse representations to be sufficient for decision-making. The agent
does not need to know everything about $E$; it only needs to know which
features of $E$ constrain its available actions and their consequences.

\section{The IAD Correspondence}

The Introspective Adversarial Distillation framework illuminates a key
feature of mature agency: the transition from external-authority-driven
to self-consistency-driven cognition.

In the IAD framework, the trust weight $\alpha = P_{\text{teacher}}
(y \mid x_{\text{adv}})^\beta$ modulates how much a student network
defers to its teacher versus its own prior beliefs. When the teacher is
reliable, $\alpha \approx 1$ and external authority dominates. When the
teacher fails on the student's own adversarial examples, $\alpha \to 0$
and self-consistency takes over.

\begin{definition}{Cognitive Authority Transition}{auth-transition}
  An agent undergoes a \emph{cognitive authority transition} when its
  self-generated model $\hat{E}_A^{\text{self}}$ becomes more reliable
  than its externally-provided model $\hat{E}_A^{\text{ext}}$ in some
  region of its observation space.

  In RSVP terms, the total flow is:
  \[
    \mathbf{v}_{\text{total}} = \alpha(x)\,\mathbf{v}_{\text{ext}}
    + (1-\alpha(x))\,\mathbf{v}_{\text{self}}
  \]
  where $\alpha(x)$ is a local trust measure over the semantic manifold.
\end{definition}

This formalizes a deep observation about cognitive development. Children
begin in a high-$\alpha$ regime — high deference to parents, teachers,
and authority. As they develop their own models of the world through
experience, $\alpha$ decreases in regions where their model outperforms
the external authority's. Expertise is the accumulation of low-$\alpha$
regions — domains where self-generated models dominate.

\section{Latent Worlds}

The most striking implication of projection-based consciousness is that
the agent's experienced world is always a latent construct: a low-dimensional
reconstruction of a high-dimensional reality, filtered through the agent's
specific projection $\Pi_A$.

\begin{definition}{Latent World}{latent-world}
  The \emph{latent world} of agent $A$ at time $t$ is the distribution
  $\mathcal{P}_A(O_t) \in \Delta(\hat{E}_A)$ over internal models
  consistent with the current observation $O_t = \Pi_A(E_t)$.
  The agent's \emph{experienced world} is the mode or expectation of
  this distribution.
\end{definition}

Two agents with different projection maps $\Pi_A \neq \Pi_B$ will
construct different latent worlds from the same environment $E$.
This is not relativism — the environment $E$ is shared and objective —
but it does imply that \emph{what is salient} differs systematically
across agents with different sensory and cognitive projections.

\begin{example}
  A bat and a human share the same physical environment. The bat's
  projection $\Pi_{\text{bat}}$ emphasizes high-frequency acoustic
  structure; the human's projection $\Pi_{\text{human}}$ emphasizes
  low-frequency visual structure. The bat's latent world contains
  rich detail about insect wing-beat frequencies; the human's contains
  rich detail about color and shape. Neither is more accurate — both
  are faithful sparse reconstructions of different projections of the
  same environment.
\end{example}

\section{Consciousness as Accessibility Navigation}

The RSVP accessibility framework (Chapters~14, 21, Appendix~J) suggests
a specific functional role for consciousness. Rather than being epiphenomenal
or merely an observer, consciousness is the agent's accessibility navigation
system: the process by which the agent moves through its latent world in
ways that preserve and expand future admissible states.

\begin{namedthm}{The Consciousness-as-Navigation Hypothesis}
  Conscious processing preferentially explores regions of the agent's
  latent world $\hat{E}_A$ that maximize the expected accessibility
  entropy of the agent's future states:
  \[
    a^* = \arg\max_{a \in \mathcal{A}(x)} \mathbb{E}\bigl[S_A(x_{t+1})
    \mid x_t = x, a_t = a\bigr]
  \]
  Consciousness is the computational substrate for this maximization.
  Attention is the selection of features of $\hat{E}_A$ most relevant
  to accessibility navigation.
\end{namedthm}

This is not a complete theory of consciousness — it says nothing about
the phenomenal or subjective character of experience. But it offers a
\emph{functional} definition that is both precise and testable: a system
is conscious to the degree that it navigates its latent world using
accessibility gradients rather than fixed reflex arcs.

\begin{exercise}
  Define a formal measure of ``consciousness'' for a projection-based
  agent as the mutual information between the agent's action sequence
  and its accessibility entropy trajectory. Show that reflex-arc agents
  (agents with fixed stimulus-response maps) have zero mutual information
  by this measure. What is the measure for an optimal accessibility-maximizing
  agent?
\end{exercise}

\begin{exercise}
  The IAD trust weight $\alpha(x)$ is locally computed — it depends on
  the teacher's reliability at the specific query $x$. Extend this to
  the cognitive authority transition: define a local trust measure
  $\alpha_A(x)$ for a human agent over the semantic manifold $M_s$.
  What observational data would allow $\alpha_A$ to be estimated from
  behavior?
\end{exercise}

\begin{exercise}[title={Open problem}]
  Pearce's zero ontology proposes that the universe has zero net
  information content: $\sum_i \psi_i = 0$. If this is true, what
  is the information content of a specific agent's latent world
  $\hat{E}_A$? Can a zero-information universe contain non-zero
  information latent worlds? (Hint: consider the relationship between
  global information content and local conditional information content
  relative to a specific projection.)
\end{exercise}
